\chapter{Informasi Kuantum dan Machine Learning Kuantum}
\label{ch:kuantum}

Pada musim panas 2025 saya mengikuti kuliah Quantum Information dari Johannes Kofler, dengan delapan tugas yang membawa kami dari
keadaan satu qubit sampai keterkaitan, teleportasi, dan pengukuran kuantum yang presisi. Naskah kuliahnya dibuka dengan kalimat
yang menjadi tema seluruh bab ini: informasi itu fisis. Ia selalu dibawa oleh benda fisik, dan kalau benda itu mengikuti hukum
mekanika kuantum, aturan pengolahan informasinya ikut berubah. Bab ini merangkum kuliah itu, lalu melihat apa yang terjadi ketika
gagasan-gagasannya bertemu dengan machine learning, dengan bantuan catatan kuliah Quantum Computing dari Richard Kueng dan buku
\citet{schuld2021}.

\section{Qubit}

Bit klasik bernilai 0 atau 1. Qubit adalah sistem dua keadaan yang bisa berada dalam superposisi,
\begin{equation}
\ket{\psi}=\alpha\ket0+\beta\ket1,\qquad |\alpha|^2+|\beta|^2=1,
\end{equation}
dengan amplitudo kompleks $\alpha$ dan $\beta$. Mengukur qubit di basis $\{\ket0,\ket1\}$ memberi hasil 0 dengan peluang $|\alpha|^2$ dan
1 dengan peluang $|\beta|^2$, aturan Born, dan setelah pengukuran keadaannya runtuh ke hasil yang terukur. Untuk
$\ket+=(\ket0+\ket1)/\sqrt2$, kedua hasil sama mungkin. Superposisi bukan sekadar ketidaktahuan tentang keadaan yang sebenarnya
sudah tertentu, dan perbedaannya terlihat lewat interferensi.

Gerbang kuantum adalah matriks uniter. Gerbang Hadamard,
$H=\tfrac1{\sqrt2}\begin{pmatrix}1&1\\1&-1\end{pmatrix}$, mengubah $\ket0$ menjadi $\ket+$ dan $\ket1$ menjadi
$\ket-=(\ket0-\ket1)/\sqrt2$. Terapkan $H$ dua kali pada $\ket0$: hasil pertama $\ket+$, hasil kedua
$\tfrac1{\sqrt2}(H\ket0+H\ket1)=\tfrac12\big[(\ket0+\ket1)+(\ket0-\ket1)\big]=\ket0$. Amplitudo untuk $\ket1$ dari dua jalur saling
menghapus, sedangkan amplitudo untuk $\ket0$ saling menguatkan. Kalau superposisi hanya berarti ``koin yang belum dilihat'', koin
acak yang dilempar dua kali tetap acak. Qubit kembali pasti ke $\ket0$. Naskah kuliah memakai interferometer Mach--Zehnder dan
percobaan celah ganda untuk menunjukkan gejala yang sama dengan foton sungguhan, termasuk teka-teki bom Elitzur--Vaidman, di
mana keberadaan sebuah bom bisa dideteksi tanpa satu foton pun menyentuhnya.

Keadaan sebuah qubit bisa digambar sebagai titik pada bola Bloch. Kutub utara adalah $\ket0$, kutub selatan $\ket1$, dan khatulistiwa
berisi superposisi sama rata dengan fase berbeda, termasuk $\ket+$ dan $\ket-$. Gerbang satu qubit adalah rotasi bola itu.

\section{Banyak qubit dan keterkaitan}

Keadaan $n$ qubit hidup di ruang berdimensi $2^n$, hasil kali tensor ruang setiap qubit. Lima puluh qubit sudah membutuhkan
$2^{50}\approx10^{15}$ amplitudo kompleks untuk ditulis lengkap, lebih dari memori komputer biasa. Inilah alasan sistem kuantum
sulit disimulasikan secara klasik, dan sekaligus alasan Feynman mengusulkan komputer kuantum untuk menyimulasikan alam.

Sebagian keadaan banyak qubit tidak bisa ditulis sebagai hasil kali keadaan masing-masing qubit. Keadaan Bell
\begin{equation}
\ket{\Phi^+}=\tfrac1{\sqrt2}\big(\ket{00}+\ket{11}\big)
\end{equation}
adalah contohnya. Kalau ia hasil kali $(a\ket0+b\ket1)(c\ket0+d\ket1)$, koefisien $\ket{01}$, yaitu $ad$, harus nol, sehingga $a$ atau $d$
nol, padahal koefisien $\ket{00}=ac$ dan $\ket{11}=bd$ sama-sama bukan nol. Kontradiksi. Keadaan seperti ini disebut \istilah{entangled}.
Mengukur qubit pertama memberi 0 atau 1 secara acak, tetapi qubit kedua selalu memberi hasil yang sama, sejauh apa pun keduanya
dipisahkan. Keadaan Bell dibuat dengan dua gerbang: Hadamard pada qubit pertama, lalu CNOT yang membalik qubit kedua kalau qubit
pertama bernilai 1.

Qubit yang terkait tidak punya keadaan sendiri. Matriks densitas tereduksi dari satu qubit $\ket{\Phi^+}$ adalah $\tfrac12\bm I$, campuran
sempurna, dengan entropi von Neumann satu bit, padahal keadaan keseluruhannya murni dengan entropi nol. Informasi tersimpan
seluruhnya di korelasi, tidak di bagian-bagiannya. Entropi von Neumann adalah perluasan entropi Shannon (\cref{lamp:matematika})
ke keadaan kuantum, dihitung dari nilai eigen matriks densitas. Interaksi yang tidak terkendali dengan lingkungan menciptakan
keterkaitan semacam ini antara qubit dan lingkungannya, sehingga qubit kehilangan superposisinya dan berperilaku seperti bit klasik
acak. Proses ini, \istilah{decoherence}, adalah musuh utama komputer kuantum.

\section{Apa yang tidak bisa dan bisa dilakukan}

Teorema \istilah{no-cloning} menyatakan bahwa keadaan kuantum yang tidak diketahui tidak bisa disalin. Pembuktiannya hanya butuh
linearitas. Misalkan ada mesin uniter yang menyalin, $U\ket\psi\ket s=\ket\psi\ket\psi$ untuk setiap $\ket\psi$. Maka ia menyalin $\ket0$
dan $\ket1$. Karena linear, ia memetakan $(\ket0+\ket1)/\sqrt2$ ke $(\ket{00}+\ket{11})/\sqrt2$, padahal salinan yang benar adalah
$\tfrac12(\ket0+\ket1)(\ket0+\ket1)$, yang memuat suku $\ket{01}$ dan $\ket{10}$. Keduanya berbeda, sehingga mesin seperti itu tidak ada.
Sifat ini membuat informasi kuantum tidak bisa disadap tanpa meninggalkan jejak, dasar dari distribusi kunci kuantum.

Keterkaitan membuka protokol yang tidak punya padanan klasik. Pada \istilah{superdense coding}, dua orang yang sudah berbagi pasangan
Bell bisa mengirim dua bit klasik dengan mengirim hanya satu qubit. Pada \istilah{teleportation}, keadaan sebuah qubit yang tidak
diketahui dipindahkan dari satu tempat ke tempat lain dengan memakai satu pasangan Bell dan mengirim dua bit klasik. Keadaan aslinya
hancur dalam proses itu, sesuai no-cloning, dan karena bit klasik harus dikirim, tidak ada informasi yang bergerak lebih cepat dari
cahaya.

\section{Bell dan dunia yang tidak lokal}

Einstein, Podolsky, dan Rosen berargumen pada 1935 bahwa korelasi sempurna pada keadaan terkait berarti hasil pengukuran sudah
ditentukan sebelumnya oleh ``variabel tersembunyi''. Bell menunjukkan bahwa gagasan itu bisa diuji. Dalam bentuk CHSH, dua pihak
masing-masing memilih satu dari dua pengaturan pengukuran, $a$ atau $a'$ dan $b$ atau $b'$, dan mencatat hasil $\pm1$. Untuk teori
apa pun yang realis dan lokal, kombinasi korelasi
\begin{equation}
S=E(a,b)+E(a,b')+E(a',b)-E(a',b')
\end{equation}
tidak bisa melebihi 2. Alasannya bisa diperiksa langsung. Kalau setiap hasil sudah tertentu sebelumnya, $S$ untuk satu kejadian
sama dengan $A(B+B')+A'(B-B')$, dan karena $B,B'\in\{\pm1\}$, salah satu dari $B+B'$ dan $B-B'$ bernilai nol dan yang lain $\pm2$,
sehingga nilainya $\pm2$. Rata-ratanya pun tidak bisa melewati 2. Mekanika kuantum dengan keadaan Bell dan sudut pengukuran yang
tepat memberi $S=2\sqrt2\approx2{,}83$. Percobaan-percobaan yang semakin bersih selama setengah abad memperlihatkan pelanggaran ini,
dan Hadiah Nobel Fisika 2022 diberikan kepada Alain Aspect, John Clauser, dan Anton Zeilinger untuk percobaan dengan foton yang
terkait. Kuliah menghubungkannya dengan distribusi kunci kuantum berbasis keterkaitan: penyadap yang mengganggu pasangan foton akan
menurunkan nilai $S$, dan penurunan itu bisa dideteksi.

\section{Algoritme kuantum}

Algoritme pertama yang dibahas kuliah adalah algoritme Deutsch. Diberikan fungsi rahasia $f:\{0,1\}\to\{0,1\}$, tentukan apakah $f$
konstan atau seimbang. Secara klasik, $f$ harus ditanya dua kali. Dengan rangkaian kuantum yang memakai Hadamard di awal dan di
akhir, satu pertanyaan sudah cukup, karena kedua nilai $f$ masuk ke fase superposisi dan interferensi di akhir mengubah beda fase itu
menjadi hasil pengukuran yang pasti. Hasilnya tidak berguna secara praktis, tetapi polanya, superposisi lalu fase lalu interferensi,
adalah resep di balik semua algoritme kuantum.

Algoritme Grover mencari satu entri di antara $N$ yang tidak terurut dengan sekitar $\tfrac\pi4\sqrt N$ pertanyaan, sedangkan pencarian
klasik butuh rata-rata $N/2$. Untuk sejuta entri, sekitar 785 pertanyaan dibanding setengah juta. Percepatannya kuadratik, bukan
eksponensial. Algoritme Shor memfaktorkan bilangan bulat jauh lebih cepat dari algoritme klasik terbaik yang diketahui, sehingga
mengancam kriptografi RSA, dan karena itulah kriptografi pascakuantum kini sedang dibakukan. Kedua algoritme membutuhkan qubit yang
jauh lebih banyak dan jauh lebih bebas galat daripada yang tersedia sekarang. Perangkat saat ini, dengan puluhan sampai ratusan qubit
yang berisik, oleh \citet{preskill2018} dinamai era NISQ, \emph{noisy intermediate-scale quantum}.

\section{Mengukur dengan presisi kuantum}

Bab terakhir naskah kuliah membahas metrologi kuantum, yang menyambung langsung ke teori estimasi di \cref{ch:sinyal}. Untuk menaksir
sebuah fase $\varphi$ dengan $N$ probe yang independen, galatnya turun seperti $1/\sqrt N$, batas yang disebut \emph{standard quantum
limit}. Inilah varians rata-rata dari \cref{lamp:matematika}, dalam bahasa fisika. Kalau $N$ probe dibuat saling terkait, fasenya
terkumpul $N$ kali lebih cepat, dan galatnya bisa turun seperti $1/N$, batas Heisenberg. Dengan sejuta foton, selisihnya seribu
kali. Batas-batas ini dirumuskan dengan informasi Fisher kuantum dan batas Cramér--Rao, alat yang sama dengan CRLB di
\cref{ch:sinyal}. Detektor gelombang gravitasi sudah memakai cahaya dengan keadaan kuantum khusus untuk menekan noise di bawah batas
klasik.

\section{Machine learning dengan komputer kuantum}

Pertemuan antara komputasi kuantum dan machine learning berjalan ke dua arah. Machine learning dipakai untuk mengendalikan dan
mengkalibrasi perangkat kuantum, membaca hasil pengukuran yang berisik, dan menyimulasikan sistem kuantum, seperti FermiNet di
\cref{ch:fiskom}. Arah sebaliknya, komputer kuantum untuk machine learning, adalah pokok buku \citet{schuld2021}.

Langkah pertama selalu memasukkan data klasik ke keadaan kuantum. \emph{Basis encoding} menulis bit-bit data sebagai keadaan basis.
\emph{Amplitude encoding} menyimpan vektor $2^n$ dimensi sebagai amplitudo $n$ qubit. Sepuluh qubit memuat vektor 1024 dimensi.
\emph{Angle encoding} memakai setiap fitur sebagai sudut rotasi satu qubit. Pilihan ini bukan detail teknis, karena ia menentukan fungsi
apa yang bisa dihitung model, dan menyiapkan keadaan dengan amplitude encoding secara umum mahal, sehingga keuntungan eksponensial
dari penyimpanan yang padat bisa hilang di langkah ini.

Model yang paling banyak diteliti untuk perangkat NISQ adalah \istilah{variational quantum circuit}: rangkaian dengan sudut rotasi yang
menjadi parameter, keluaran berupa nilai harapan sebuah pengukuran, dan parameter yang dioptimasi oleh komputer klasik. Gradien
dihitung dengan \istilah{parameter-shift rule}. Untuk gerbang rotasi $e^{-i\theta P/2}$ dengan $P$ matriks Pauli, nilai harapan keluaran
adalah fungsi sinusoidal dari $\theta$, $f(\theta)=A+B\cos\theta+C\sin\theta$, dan turunannya bisa dihitung tepat dari dua evaluasi
rangkaian yang sama,
\begin{equation}
\frac{\partial f}{\partial\theta}=\frac{f(\theta+\tfrac\pi2)-f(\theta-\tfrac\pi2)}2 .
\end{equation}
Memeriksanya mudah: $\cos(\theta\pm\pi/2)=\mp\sin\theta$ dan $\sin(\theta\pm\pi/2)=\pm\cos\theta$, sehingga selisihnya dibagi dua adalah
$-B\sin\theta+C\cos\theta$, tepat turunan $f$. Berbeda dengan beda hingga, aturan ini eksak, bukan aproksimasi.

Cara pandang kedua menghubungkan semua ini dengan \cref{ch:teori}. Peta fitur kuantum $\vx\mapsto\ket{\phi(\vx)}$ mendefinisikan kernel
$k(\vx,\vx')=|\braket{\phi(\vx)}{\phi(\vx')}|^2$, yang bisa ditaksir dengan mengukur tumpang tindih dua keadaan
\citep{havlicek2019}. Model kuantum yang dilatih dengan cara ini adalah metode kernel, sehingga teorema representer, regularisasi,
dan analisis generalisasinya berlaku. Pertanyaannya bergeser menjadi: apakah ada kernel yang berguna untuk data nyata, mudah dihitung
dengan komputer kuantum, dan sulit dihitung dengan komputer klasik?

Kuliah dan bacaannya juga jujur tentang hambatannya. Rangkaian variasional yang acak dan dalam punya lanskap loss yang hampir datar di
mana-mana, \emph{barren plateau} \citep{mcclean2018}: gradiennya mengecil eksponensial terhadap jumlah qubit, masalah vanishing gradient
dari \cref{ch:lstm} dalam bentuk yang lebih parah. Noise perangkat menambah galat pada setiap evaluasi. Memuat data klasik bisa
menghabiskan keuntungan apa pun. Dan untuk banyak tugas, belum ada bukti bahwa model kuantum lebih baik dari model klasik yang
dirancang dengan baik. Tempat yang paling menjanjikan untuk keunggulan kuantum barangkali bukan data klasik seperti gambar atau teks,
melainkan data yang memang kuantum sejak awal, dari sensor kuantum atau simulasi sistem kuantum.

\section{Perkembangan riset}

Perangkat kuantum berkembang cepat, dengan perhatian utama pada koreksi galat: banyak qubit fisik yang berisik digabung menjadi
satu qubit logis yang lebih andal, dan percobaan-percobaan terbaru mulai memperlihatkan bahwa galat logis bisa turun ketika ukuran
kode koreksi dinaikkan. Machine learning ikut berperan di sini, misalnya untuk men-decode sindrom galat dengan cepat. Pada saat yang
sama, banyak klaim keunggulan kuantum untuk tugas tertentu kemudian disusul oleh algoritme klasik yang lebih baik, sehingga bidang
ini sedang belajar membedakan percepatan yang mendasar dari percepatan yang hanya mencerminkan algoritme klasik yang belum optimal.
Sikap yang sehat, seperti di bagian-bagian lain buku ini, adalah membandingkan dengan pembanding klasik terbaik, bukan yang
termudah.

\sumberbab{Bab ini mengikuti naskah kuliah Quantum Information (keadaan dan pengukuran kuantum, percobaan satu qubit, keadaan banyak
qubit dan gerbang, no-cloning, superdense coding, teleportasi, ketaksamaan Bell dan distribusi kunci kuantum, keadaan campuran dan
dekoherensi, entropi, kanal kuantum, metrologi kuantum). Bacaan pendamping: catatan kuliah Introduction to Quantum Computing,
\citet{nielsen2010}, dan untuk machine learning kuantum \citet{schuld2021}.}
